\clearpage
\section{Synchronization from averaging for a weakly forced oscillator}
\subsection{Main synchronization region}
\subsubsection{slow phase}
Since we define 
\begin{equation}
    \psi = \theta - \Omega t
\end{equation}
then 
\begin{equation}
    \dot{\psi}=\dot{\theta}-\Omega=\Delta \omega+\epsilon Q(\psi+\Omega t, \Omega t)
\end{equation}
both of these terms are \(O(\epsilon)\) so 
\begin{equation}
    \dot{\psi}=O(\epsilon)
\end{equation}
so over one period 
\begin{equation}
    \Delta\psi = O(\epsilon T)
\end{equation}
meaning \(\psi\) is constant to leading order during a single period.
\subsubsection{Averaged equation}
inserting theta and phi
\begin{equation}
    Q(\psi+\Omega t,\Omega t)=\sum_{m,n}c_{m,n}e^{im\psi}e^{i\ins{m+n}\Omega t}
\end{equation}
averaging over one period means that only the terms for which \(m+n=0\) survive which eliminates the dependence on \(\Omega\), meaning 
\begin{equation}
    \dot{\psi}=\Delta\omega + \epsilon f(\psi)
\end{equation}
where 
\begin{equation}
    f(\psi)=\sum_{m}c_{m,-m}e^{im\psi}
\end{equation}
\subsubsection{locking range}
Since a locked solution is a fixed point of the averaged relative phase equation 
\begin{equation}
    0=\Delta\omega + \epsilon f(\psi^*)
\end{equation}
locking happens when 
\begin{equation}
    f_{\min} \leq -\frac{\Delta\omega}{\epsilon}\leq f_{\max}
\end{equation}
so 
\begin{equation}
    -\epsilon f_{\max}\leq \Delta\omega \leq -\epsilon f_{\min}
\end{equation}
so the locking width is 
\begin{equation}
    \Delta\omega_{\text{width}} = \epsilon\ins{f_{\max}-f_{\min}}
\end{equation}
and the fixed point is stable when 
\begin{equation}
    \epsilon f'(\psi^*)<0
\end{equation}
\subsection{p:q resonances}
defining 
\begin{equation}
    \psi = q\theta-p\Omega t
\end{equation}
and assuming 
\begin{equation}
    q\omega_0 \approx p\Omega
\end{equation}
means 
\begin{equation}
    \dot{\psi}=q\dot{\theta}-p\Omega=\delta_{pq}+\epsilon qQ(\theta,\Omega t)
\end{equation}
if we define \(\phi=\Omega t\) then 
\begin{equation}
    \theta = \frac{\psi+p\phi}{q}
\end{equation}
which means 
\begin{equation}
    \dot{\psi}=\delta_{pq}+\epsilon q\sum_{m,n}c_{m,n}e^{im\psi/q}e^{i\ins{mp/q+n}\phi}
\end{equation}
and again when averaging the only terms that survive are those where 
\begin{equation}
    \frac{mp}{q}+n=0 \implies mp+nq=0
\end{equation}
but since p and q are coprime
\begin{equation}
    m=kq,\qquad n=-kp
\end{equation}
so 
\begin{equation}
    \dot{\psi}=\delta_{pq}+\epsilon f_{pq}(\psi)
\end{equation}
where 
\begin{equation}
    f_{pq}(\psi)=q\sum_k c_{kq,-kp}e^{ik\psi}
\end{equation}
\subsection{The Adler equation}
The Adler model is 
\begin{equation}
    \dot{\theta}=\omega_0-\epsilon\sin\ins{\theta-\Omega t}
\end{equation}
so
\begin{equation}
    Q(\theta,\phi) = -\sin\ins{\theta-\phi}
\end{equation}
which if we expand this in a fourier expansion, it'll contain only the harmonics
\begin{equation}
    \ins{m,n}=(1,1),\quad(-1,1)
\end{equation}
but for a p:q resonance the coefficients need the form 
\begin{equation}
    \ins{m,n}=\ins{kq,-kp}
\end{equation}
which are satisfied for the Adler harmonics only for 
\begin{equation}
    p=q=1
\end{equation}
as a result, for every \(\ins{p,q}\neq(1,1)\)
\begin{equation}
    f_{pq}=0
\end{equation}
making the averaged equation 
\begin{equation}
    \dot{\psi}=\delta_{pq}
\end{equation}
meaning there is no finite locking interval, meaning the adler equation has no finite-width higher p:q resonances.